% ---------------------------------------------------------------------------
% II. Related Work -- SIX-FAMILY version, 2026-08-29. Rewritten from
% docs/literature_review.md to match the chosen six-family scope. The earlier
% three-method version is kept as relatedwork_3method_backup.tex.
%
% Organised by the role each family plays: negative controls (foveation,
% action repeat), the recent baseline (VLA-Cache), and the conditional
% candidates (depth, guarded reuse, temporal fusion). This is the same order
% Methods (Section III) and Results walk through.
%
% ⚠️ FOUR NEW CITATION KEYS used below that are NOT yet in main.bib. Build
% proper entries from the arXiv pages before compiling (ids from
% docs/literature_review.md):
%   flashvla        2505.21200
%   ttfvla          2508.19257   (AAAI 2026)
%   vlainfoentropy  2604.05323
%   crossmodalreuse  anonymous OpenReview R6d86jMO74
% (teamvla 2512.09927 and actioncache 2607.06370 are in literature_review.md
%  but not cited in this draft; add them only if a later revision needs them.)
% Existing keys reused: openvla, shortgpt, efficientvla, gromov, molevla,
%   vlacache, vlapruner, fastv, sparsevlm, tome, schwartz, traver, gazereg,
%   lookfocusact, simplerenv, libero, vlaeval, starvla, specprune, vlaiap,
%   dqn, act, diffusionpolicy, effvlasurvey1, bagoftricks_cnn, bagoftricks_llm.
%
% Style: no em-dashes, no semicolons, no prose colons.
% ---------------------------------------------------------------------------

\section{Related Work}

We cover the cost of VLA inference, the interventions we evaluate, grouped by
the role each plays, and how the literature tests them.

\textbf{Inference cost in VLA policies.} VLA policies adapt pretrained
vision-language models to output robot actions, inheriting their size and
latency, and the resulting efficiency literature already has surveys of its
own~\cite{effvlasurvey1}. Methods divide by the resource each spends, when the
policy runs, what it is shown, and how much network each call uses. A claim
about efficiency is a claim about one of these resources being spent
differently, so we treat them as axes rather than competing methods.

\textbf{Generic shortcuts, used as negative controls.} The simplest way to
spend less is to degrade the input or to skip feedback unconditionally, and both
fail in ways that motivate the safeguards below. On the spatial side, foveation
descends from Schwartz's log-polar mapping~\cite{schwartz}, taken up in robot
vision to cut data while preserving central resolution~\cite{traver} and applied
to VLAs by gaze-conditioned policies~\cite{gazereg, lookfocusact}. At a fixed
output resolution the encoder still produces the same visual tokens, so a
pixel-space edit removes no model computation. It only tests whether the policy
survives clutter.
Methods that instead foveate inside the encoder shed tokens but give up the
pretrained weights fitted to a uniform grid~\cite{lookfocusact}. On the temporal
side, executing one decision over several control steps has a lineage from frame
skip in reinforcement learning to action chunking~\cite{dqn, act,
diffusionpolicy}. Unconditional repetition reduces model calls but acts open
loop through contacts, and training-free action reuse of this kind reports a
success decrease~\cite{flashvla} unless it separates coarse motion from
precision phases~\cite{specprune}. We therefore use fixed foveation and fixed
repeat as controls that establish why the candidate methods need gates.

\textbf{Conditional candidates.} Three families remove computation only where a
signal says it is safe. Decoder-layer redundancy is well established in language
models. But the recipes built on it disagree on which layers may go. ShortGPT
ranks every layer by Block Influence, the criterion we adopt, and constrains
nothing further~\cite{shortgpt}, while others keep the final layer and remove a
contiguous deep block~\cite{gromov}. EfficientVLA carries the unconstrained form
into VLAs without training~\cite{efficientvla}, and MoLe-VLA needs a learned
router~\cite{molevla}, placing it outside a training-free study. Guarded action
reuse makes the temporal shortcut conditional on visual and trajectory
stability, a direction that recent work supports by tying reuse to interaction
phase~\cite{flashvla, specprune}. Task-aware visual reuse is the third family.
Temporal token fusion reuses parts of the visual representation across frames to
stabilise it~\cite{ttfvla}, while selective cache reuse budgets computation by
image entropy and task attention~\cite{vlainfoentropy, vlaiap}. A shared mask
that governs both denoising and cache reuse is close to a concurrent
proposal~\cite{crossmodalreuse}, so the defensible differences are contact-aware
fallbacks, one mask driving both paths, and paired evaluation against optimized
dense inference.

\textbf{The recent baseline.} VLA-Cache reuses the key-value computation of
static patches while recomputing dynamic and task-relevant ones, retaining
history unlike token pruning~\cite{vlacache}. On base OpenVLA across LIBERO it
reports nearly maintained rather than improved average success at lower latency,
which makes it the recent anchor against which candidates are measured rather
than a first-party contribution. Visual-token methods developed for
vision-language models transfer poorly to VLAs~\cite{fastv, sparsevlm, tome},
which VLA-Cache and VLA-Pruner attribute to VLAs' short action
sequences~\cite{vlacache, vlapruner}.

\textbf{How these claims are evaluated.} The field reports results on
SimplerEnv~\cite{simplerenv} and LIBERO~\cite{libero}, and recent work addresses
the infrastructure around them, with the vla-eval harness unifying benchmarks
and documenting evaluation pitfalls earlier work had left
unrecorded~\cite{vlaeval} and StarVLA describing the field as fragmented across
incompatible codebases~\cite{starvla}. But infrastructure cannot supply the
comparison itself. Papers using several backbones change the benchmark at the
same time~\cite{vlacache, specprune, vlaiap}, those with both factors present
leave the crossing cell empty~\cite{gazereg, vlapruner}, and the tables we cite
report mean success rates rather than matched-episode outcomes. Many also
compare against an artificially slow eager baseline, which inflates a speedup.
Isolating the effect of choices that appear only in source code is an
established practice~\cite{bagoftricks_cnn, bagoftricks_llm}. We evaluate the six
families under one protocol, on matched episodes against optimized dense
inference, so that a candidate is called positive only when its speed and its
success both clear a preregistered gate.
